%%========================================================================
%% LaTeX sjabloon voor stage/projectrapport of bachelorproef
%%  HoGent Bedrijf en Organisatie
%%========================================================================

%%========================================================================
%% Preamble
%%========================================================================

\documentclass[pdftex,a4paper,12pt,twoside]{report}

% XXX: Let op: dit sjabloon is gemaakt om dubbelzijdig af te drukken
% Voor enkelzijdig, verwijder ``twoside'' hierboven.

%%---------- Extra functionaliteit ---------------------------------------

\usepackage[utf8]{inputenc}  % Accenten gebruiken in tekst (vb. é ipv \'e)
\usepackage{amsfonts}        % AMS math packages: extra wiskundige
\usepackage{amsmath}         %   symbolen (o.a. getallen-
\usepackage{amssymb}         %   verzamelingen N, R, Z, Q, etc.)
\usepackage[dutch]{babel}    % Taalinstellingen: woordsplitsingen,
                             %  commando's voor speciale karakters
                             %  ("dutch" voor NL)
\usepackage{eurosym}         % Euro-symbool €
\usepackage{geometry}
\usepackage{graphicx}        % Invoegen van tekeningen
\usepackage[pdftex,bookmarks=true]{hyperref}
                             % PDF krijgt klikbare links & verwijzingen,
                             %  inhoudstafel
\usepackage{listings}        % Broncode mooi opmaken
\usepackage{multirow}        % Tekst over verschillende cellen in tabellen
\usepackage{rotating}        % Tabellen en figuren roteren
\usepackage{natbib}          % Betere bibliografiestijlen
\usepackage{fancyhdr}        % Pagina-opmaak met hoofd- en voettekst

\usepackage[T1]{fontenc}     % Ivm lettertypes
\usepackage{lmodern}
\usepackage{textcomp}

\usepackage{lipsum}          % Voor vultekst (lorem ipsum)

\usepackage{booktabs}
%%---------- Layout ------------------------------------------------------

% hoofdingen, enz.
\pagestyle{fancy}
% enkel hoofdstuktitel in hoofding, geen sectietitel (vermijd overlap)
\renewcommand{\sectionmark}[1]{}

% lijn, wordt gebruikt in titelpagina
\newcommand{\HRule}{\rule{\linewidth}{0.5mm}}

% Leeg blad
\newcommand{\emptypage}{
\newpage
\thispagestyle{empty}
\mbox{}
\newpage
}

% Gebruik een schreefloos lettertype ipv het "oubollig" uitziende
% Computer Modern
\renewcommand{\familydefault}{\sfdefault}

% Commando voor invoegen Java-broncodebestanden (dank aan Niels Corneille)
% Gebruik: \codefragment{source/MijnKlasse.java}{Uitleg bij de code}
\newcommand{\codefragment}[2]{ \lstset{%
  language=java,
  breaklines=true,
  float=th,
  caption={#2},
  basicstyle=\scriptsize,
  frame=single,
  extendedchars=\true
}
\lstinputlisting{#1}}

% Opmaak van de paragrafen
\setlength{\parskip}{1em}

%%---------- Documenteigenschappen ---------------------------------------
%% Vul dit aan met je eigen info:

% Je eigen naam
\newcommand{\student}{Tomas Vercautter}

% De naam van je lector, begeleider, promotor
\newcommand{\promotor}{Bert Van Vreckem}

% De naam van je co-promotor
\newcommand{\copromotor}{Bert Van Vreckem}

% Indien je bachelorproef in opdracht van een bedrijf of organisatie
% geschreven is, geef je hier de naam.
\newcommand{\instelling}{---}

% De titel van het rapport/bachelorproef
\newcommand{\titel}{Security en Managebility van Docker containers}

% Datum van indienen
\newcommand{\datum}{27 mei 2016}

% Faculteit
\newcommand{\faculteit}{Faculteit Bedrijf en Organisatie}

% Soort rapport
\newcommand{\rapporttype}{Scriptie voorgedragen tot het bekomen van de graad van\\Bachelor in de toegepaste informatica}

% Academiejaar
\newcommand{\academiejaar}{2015-2016}

% Examenperiode
%  - 1e semester = 1e examenperiode
%  - 2e semester = 2e examenperiode
%  - tweede zit = 3e examenperiode
\newcommand{\examenperiode}{Tweede examenperiode}

%%========================================================================
%% Inhoud document
%%========================================================================

\begin{document}

%%---------- Front matter ------------------------------------------------
%% Het voorblad - Hier moet je in principe niets wijzigen.

\begin{titlepage}
  \newgeometry{top=2cm,bottom=1.5cm,left=1.5cm,right=1.5cm}
  \begin{center}

    \begingroup
    \rmfamily
    \includegraphics[width=2.5cm]{img/HG-beeldmerk-woordmerk}\\[.5cm]
    \faculteit\\[3cm]
    \titel
    \vfill
    \student\\[3.5cm]
    \rapporttype\\[2cm]
    Promotor:\\
    \promotor\\
    Co-promotor:\\
    \copromotor\\[2.5cm]
    Instelling: \instelling\\[.5cm]
    Academiejaar: \academiejaar\\[.5cm]
    \examenperiode
    \endgroup

  \end{center}
  \restoregeometry
\end{titlepage}

% Schutblad

\emptypage


\begin{titlepage}
  \newgeometry{top=5.35cm,bottom=1.5cm,left=1.5cm,right=1.5cm}
  \begin{center}

    \begingroup
    \rmfamily
    \faculteit\\[3cm]
    \titel
    \vfill
    \student\\[3.5cm]
    \rapporttype\\[2cm]
    Promotor:\\
    \promotor\\
    Co-promotor:\\
    \copromotor\\[2.5cm]
    Instelling: \instelling\\[.5cm]
    Academiejaar: \academiejaar\\[.5cm]
    \examenperiode
    \endgroup

  \end{center}
  \restoregeometry
\end{titlepage}


\begin{abstract}
% TODO: De "abstract" of samenvatting is een kernachtige (max 1 blz. voor een
% thesis) synthese van het document. In ons geval beschrijf je kort de
% probleemstelling en de context, de onderzoeksvragen, de aanpak en de
% resultaten.
\end{abstract}

\chapter*{Voorwoord}
\label{ch:voorwoord}

% TODO: Vergeet ook niet te bedanken wie je geholpen/gesteund/... heeft


\tableofcontents

% Als je een lijst van afkortingen of termen wil toevoegen, dan hoort die
% hier thuis. Gebruik bijvoorbeeld de ``glossaries'' package.

%%---------- Kern --------------------------------------------------------


\chapter{Inleiding}
\label{ch:inleiding}

Het gebruik van docker containers is een manier van virtualiseren van services. Hierbij worden services in een 

In deze bachelorproef zal ik onderzoek doen naar de security en managebility van Docker containers en deze vergelijken met de security en managebility van het draaien van de services rechtstreeks op de virtuele machine.

Deze bachelorproef is onderverdeeld in drie grote delen. In de eerste plaats bespreek ik de security zelf. In een tweede deel zal ik het hebben over de managebility van Docker containers. En als laatste onderdeel bespreek ik logging en monitoring bij Docker containers.

\section{Probleemstelling en context}
\label{sec:onderzoeksvragen}

% TODO: Wees zo concreet mogelijk bij het formuleren van je
% onderzoeksvra(a)g(en). Een onderzoeksvraag is trouwens iets waar nog
% niemand op dit moment een antwoord heeft (voor zover je kan nagaan).

Docker is een nieuwe speler op de virtualisatiemarkt die in korte tijd een substantieel marktaandeel heeft ingenomen. Hiervoor hebben ze Linux LXC containers gebruiksvriendelijker gemaakt. Dit type container ligt ergens in het midden tussen een chroot jail (isoleren van een proces van het rest van het systeem) en een volledige virtuele machine. Sinds Docker is uitgekomen voor het grote publiek zijn er altijd al vragen geweest naar in hoeverre deze vorm van virtualisatie veilig genoeg was om in productie te gebruiken. Als we in Google iets zoeken rond problemen met Docker of mensen die gestopt zijn met Docker te gebruiken, vinden we heel wat posts met als titel iets in de zin van "Why I stopped using Docker". De meeste mensen stoppen met het gebruik van Docker omdat ze problemen vinden met ofwel de beveiliging van de container omgeving, ofwel hebben ze problemen met het managen van hun containers. Deze kritieken of problemen met Docker zijn er al sinds de eerste uitgave van Docker en blijven nu en dan de kop op steken.  

\section{Doelstelling en onderzoeksvraag/-vragen}

De doelstelling van deze bachelor proef bestaat uit het ophelderen van kritieken op vlak van de managebility en beveiliging van deze containers. Hiervoor gaan we eerst op zoek naar welke kritiekpunten er te vinden zijn. er zal zo veel mogelijk gekijken worden naar recentere kritieken omdat Docker snel geëvolueerd is, kunnen kritieken van twee maanden geleden al verouderd zijn. Daarna zullen we in de eerste plaats kijken of deze kritiek gerechtvaardigd is. Indien dit het geval is, zal er onderzocht worden hoe we deze kritiekpunten mogelijks kunnen wegwerken met bepaalde zelf toepasbare technieken of best-practices. Of zijn er elementen die zullen moeten geïmplementeerd worden door de Docker community zelf?

Daarnaast zijn er veel mensen die van Docker containers een veiligere en meer handelbare ervaring verwachten dan als ze met virtuele machines werken. Dus zal er in dit onderzoek ook gekeken worden naar hoe Docker containers op minstens hetzelfde niveau gekregen kunnen worden als een klassieke virtuele machine, zowel qua beheersbaarheid als op beveiligings vlak

\include{Methodologie}
\chapter{Docker}
\section{Geschiedenis van Containers}
%TODO bron docker relases invoegen
%TODO chroot bron: http://www.faqs.org/faqs/unix-faq/faq/part6/section-2.html
%TODO what chroot does: http://man7.org/linux/man-pages/man2/chroot.2.html
Docker op zich bestaat nog niet zo lang. Versie 0.1.0 is uitgekomen op 26 maart 2016 en versie 1.0.0 en dus de eerste effectieve release kwam uit op 9 juni 2014. Hierdoor is er nog niet veel geschiedenis aan Docker. Maar als we terugkijken naar de evolutie van containers gaat dit verder terug. 
Het idee van containers gaat terug tot 1979 met de toevoeging van chroot in UNIX V7. chroot is een eerste concept van containerization. Hierbij wordt voor elk proces een geisoleerde ruimte op de schijf voorzien, ook wel een "chroot jail" genoemd.
%TODO http://imesh.io/evolution-of-linux-containers-and-future/
Na chroot duurde het tot het jaar 2000 tot het volgende op de markt kwam. FreeBSD Jails(2000) en Linux VServer(2001) waren beide jail mechanismen die verder bouwden op chroot. Hiernaa kwamen enkele oplossingen op de markt die al dichter bij de container die we nu kennen aansluiten namelijk Solaris Containers(2004), OpenVZ(2005), Process Containers(2006) en Control Groups(2007).

Daarna werd in 2008 LXC (LinuX Containers) toegevoegd aan de linux kernel. Dit was de basis voor latere containerizatie technologieën. LXC cobineerde het gebruik van cgroups (dit stond in voor isolatie en resource management) en namespaces (dit zorgde ervoor dat groups opgesplitst konden worden waardoor ze elkaar niet kunnen "zien")\cite{The Docker Ecosystem: An Overview of Containerization | DigitalOcean2015}. De eerste versies van Docker waren ontwikkeld als een interne tool bij het bedrijf dotCloud. Deze tool zorgde voor een versimpeling voor het gebruik van LXC. Later werd dit uitgebracht onder de naam Docker een 	en werd de LCX container vervangen door hun eigen driver LibContainer. Docker zorgde als een van de eerste voor een volledig ecosysteem voor het managen van containers\cite{evolution of linux containers & future2016}.
 
\section{Wat is Docker}

Docker is een open source project om een applicatie te verpakken, verzenden en te draaien als een weinig resources gebruikende container \cite{Docker:TheContainerEngine2016}. Anders gezegd word in een Docker Container alles om een programma te draaien verpackt in een compleet bestandsysteem namelijk: code, runtime, stysteem tool en systeem librairies. Alles wat je normaal op een server zou installeren. Dit garandeerd dat de applicatie zich altijd op dezelfde manier zal gedragen in welke omgeving het ook draait \citep{WatIsDocker2016}.

Een vergelijking met virtuele machines is rap gemaakt. Voor resource isolatie werkt het op dezelfde manier. Maar waar de verschillen liggen is de grootte, flexibiliteit en performatie. Doordat Docker gebruik maakt van containerization gebeurt de virtualizatie op de kernel niveau tegenover virtuele machines waarbij dit niet zo is. 

Een korte visualisatie van hoe containers werken kunnen we zien op de volgende foto's. Hierbij zien we dat traditioneel bij virtuele machines we op de host die zich rechtstreeks bevind op de hardware een hypervisor nodig hebben die onze hardware gaat virtualiseren voor gebruik van de virtuele gast machines. Op elke gast machine bevind zich een eigen besturingssysteem die uniek is voor elke gast. Daarnaast moeten ook onze applicatie draaien samen met de binaries en librairies die het nodig heeft 

\section{Hoe gebruikt men Docker}

\chapter{Security}

Op vlak van security ben ik gestart van CIS Docker 1.6 Benchmark geschreven door \cite{PravinGoyal2015}. Hierin is in samenwerking met de "Center for Internet Security" onderzoek gedaan naar de beveiliging van docker containers. Dit onderzoek is onderverdeeld in zes grote onderdelen namelijk Host configuration, Docker daemon configuration, docker daemon configuration files, Container images and build file, Container runtime en docker security operations. Hierbij wordt er gekeken welke van deze elementen standaard in docker worden uitgevoerd. In dit hoofdstuk zal ik dezelfde zes onderdelen behandelen en vergelijen. Hierbij zal ik kijken hoe we deze beveiliging aanpakken wanneer we onze services in een Docker container draaien tegenover wanneer we deze native in onze virtuele machine draaien.


Hierbij heb ik ook telkens deze methodes voor Docker containers uitgetest om te kijken hoe makkelijk deze kunnen geïmplementeerd worden in een nieuwe omgevinge of bestaande omgeving. Een grote hulp hiervoor was de github repository "docker-bench-security" van docker een grote hulp. Dit is een script die alle elementen die worden aangehaald in CIS Docker 1.6 Benchmark automatisch getest en feedback over gegeven.

\section{Host Configuration}

Een van de belangrijkste onderdelen in het beveiligen van de Docker stack is het beveiligen van de host machine waarop de docker containers zich zullen bevinden. Dit heeft drie onderdelen namelijk het up to date houden van het systeem, het juist configureren van het systeem, en het auditten van het systeem.

\subsection{Up to date houden van het systeem}

Om de veiligheid van een systeem te kunnen garanderen moet er ten allen tijde gezorgd worden dat het systeem up to date blijft. Dit zorgt ervoor dat security flaws die gekend en gepatcht zijn ook op het systeem beveiligd zijn. Doordat we Docker gebruiken komt hier nog een extra laag bij. Bij het Updaten van een machine waarop de services rechtstreeks draaien moet enkel het host systeem, de kernel van het host systeem en de services up to date blijven. Doordat Docker een kernel deelt met het hostsysteem wordt het up to date houden van de kernel belangrijker. Naast het onderhouden van de voorvernoemde elementen moet er nu, doordat we Docker gebruiken, nog met twee extra dingen rekening gehouden worden. Enerzijds moet Docker zelf nu ook upgedate worden, maar een onderdeel die soms vergeten wordt is dat elke container nu ook zijn eigen besturingssysteem gebruikt. Hierover wordt er verder bij het beveiligen van Container images meer uitleg over gegeven.

% TODO dev tools? and explanaition about stable releases

\subsection{Configureren van het systeem}

Bij het configureren van het hostysteem wordt net zoals bij het up to date houden van het systeem dezelfde stappen van een native applicatie op een hostysteem configureren gevolgd. Maar door de extra laag komen daar nog extra stappen bij. Bij een systeem waar Docker op draait wordt er aangeraden om een aparte partitie te creëren voor de containers. Dit is een simpele operatie waardoor alle docker gerelateerde files zich niet meer tussen de files van het hostsysteem bevinden.

Daarnaast moet ook onder controle gehouden worden wie toegang krijgt tot de docker daemon. De Docker daemon heeft 'root' access, gebruikers die worden toegevoegd aan de 'docker' usergroup en geen 'root' privileges hebben kunnen hierdoor 'root' access verkrijgen tot het hostsysteem. Men directories sharen tussen het host systeem en een guest container. Doordat containers standaard altijd runnen als 'root' heeft de container als de '/' directory gemount is op deze container ongelimiteerde toegang tot het volledige host systeem. Hierdoor kunnen zonder ristricties aanpassingen aan het hostysteem gedaan worden. Dit betekend dat een gebruiker die toegang heeft tot de docker daemon verhoogde rechten kan verkrijgen door gewoon een container op te starten.

Zoals bij een virtuele machine waarop de services rechtstreeks wordt ook bij een systeem waar docker containers op draaien aangeraden overbodige services uit te schakelen of te verwijderen. Moest een gebruiker door een van deze services toegang krijgen tot het hostsysteem. Zou hij makkelijker 'root' rechten kunnen vergrijgen volgens de manier beschreven in de vorige paragraaf.

\subsection{Auditen van het systeem}

% TODO find paper on auditing systems

\chapter{Managebility}

Wanneer de managebility van Docker containers onderzocht wordt, zijn er drie opties, in de eerste plaats kan alles manueel gemanaged worden. Daarnaast kan  ook alles automatisch laten gemanaged worden door een proces manager. Ofwel kunnen er tools gebruikt worden om onze containers te beheren. Als men bezig is met het beheren van Docker containers moet er ook rekening gehouden worden met het beheren van onze docker images. 

\section{Manueel beheren}

Een optie waarmee meestal begonnen wordt is het manueel managen van de Docker containers. Bij deze optie zullen de docker containers beheren enkel met de commando's die docker ons aanbied. Elke container invividueel moeten aanroepen in elk commando lijkt mogelijk in het begin. Maar na het manueel uitvoeren van deze commando's wordt het al snel duidelijk datmet behulp van een kort script of een makefile het beheren heel wat makkelijker kan. Op deze manier kunnen we zowel onze containers managen als onze Docker images builden met veel kortere commando's. Containers managen met de volledige docker commando's of scripts/makefiles is doenbaar indien we enkele containers moeten managen zoals in een ontwikkel omgeving of een stabiele kleine omgeving. Het manueel managen van onze containers verwacht ook meer stabiliteit van onze containers. Indien er niemand de containers aan het monitoren zou zijn wanneer een container crasht, moet er gewacht worden tot er iemand deze manueel terug opstart. Hierdoor kan het dus gebeuren dat een container relatief langer onbeschikbaar is. In een development omgeving zijn de vereisten voor uptime veel lager. Hier is het manueel managen van onze containers en images wel een optie. Bij het manueel managen wordt er ook gebruik gemaakt van scripts of makefiles. Een voorbeeld van beide vind u hieronder.

\begin{lstlisting}[language=bash, style=configstyle]
CONTAINER ?= redis
IMAGE_NAME ?= redis
IMAGE_VERSION ?= latest
HOSTNAME ?= redis-host


# full image name
IMAGE = $(IMAGE_NAME):$(IMAGE_VERSION)
# full ssh connect vars
CONNECT = $(REMOTE_USER_NAME)@$(IP)

build:
	docker build -t $(IMAGE) .

run:
	docker run -d --name $(CONTAINER) -h $(HOSTNAME) $(IMAGE)

logs:
	docker logs $(CONTAINER)

kill_container:
	docker kill $(CONTAINER)

remove_container:
	docker rm $(CONTAINER)

remove_image:
	docker rmi $(IMAGE)

update: kill_container remove_container remove_image build run
\end{lstlisting}

Met deze makefile hebben wordt er toegang verkregen tot de verschillende commando's die normaal volledig met de hand zouden moeten uitgeschreven worden, maar nu kunnen deze aangeroepen worden met 'make COMMANDO'. Als we dit bekijken met de voorgaande makefile zien we dat om onze redis container uit te voeren enkel 'make run' moest ingeven worden in het terminal venster terwijl anders 'docker run --name redis redis' moest worden gebruikt. Bij dit voorbeeld is het commando dat zonder makefile zouden moeten uitgevoerd worden nog beperkt maar als er meer opties meegeven moeten worden met het commando kan het interessanter zijn  om een kort make commando uit te voeren dan elke keer dat dezelfde container  opgestart moet worden opnieuw het volledige commando uit te typen. Ook wordt consistentie behouden bij het uitvoeren van de commando's via Makefiles en kunnen eventuele menselijken fouten zoals typfouten makkelijker geëlimineerd worden. Het ander voordeel dat hier ook uit voorkomt is dat make commando's ook kunnen aanroepen worden in andere make commando's, hierdoor kan het commando 'make update' gemaakt worden. Deze verwijdert onze oude bestaande container en image. Daarna wordt een nieuwe image gemaakt en ingeladen in een nieuwe container. Bovenaan de makefile vinden we ook variabelen die gebruikt kunnen worden in onze make commando's. Met behulp van deze variabelen kunnen we een makefile snel aanpassen voor het gebruik met een ander commando.

%TODO insert script here

Bij het gebruik van scripts wordt ongeveer dezelfde functionaliteit verkregen als bij makefiles, enkel is de syntaxis voor het uitvoeren anders en de manier waarop het geschreven is. Hierboven vinden we een script die dezelfde functionaliteit zou hebben als het 'make update' commando.



\section{Managen met process manager}

Een andere optie die er is, zijn de containers managen met een process-manager. Hierbij kunnen de containers bekeken worden alsof het processen zijn. Als de  containers gemanaged worden met een process-manager, moeten configuratie bestanden handmatig geschreven worden voor de gekozen process-manager. Die gekozen manager zal dan deze bestanden inlezen en vervolgens de met de parameters vanuit de configuratie bestanden onze containers starten en managen zoals het een ander proces zou beheren. Een voorbeeld hiervan is het gebruik van Systemd Unit files. Dit is hieronder uitgewerkt voor de Redis Container uit hoofdstuk drie.

%TODO systemd unit files uitleggen

\begin{lstlisting}[language=bash, style=configstyle]
[Unit]
Description=Run a redis container
Author=Toon Lamberigts & Tomas Vercautter
Requires=docker.service

[Service]
ExecStartPre=-/usr/bin/docker rm redis-container
ExecStart=/usr/bin/docker run --rm --name redis-container -h redis-host redis
ExecStop=/usr/bin/docker stop -t 2 redis
\end{lstlisting}

Deze unit file zal bij het proberen starten van de bijbehorende service ingeladen worden. Het 'ExecStartPre' gedeelte wordt uitgevoerd voor de service effectief zal starten. Deze zal de container verwijderen indien deze bestaat, om zeker te zijn dat er geen container genaamd 'redis-container' bestaat. Dit zou een fout genereren bij het opstarten waardoor de service zou falen. Dus uit voorzorg proberen we deze container eerst te verwijderen voor het opstarten. Door het meegeven van een '-' voor het commando, zal het opstarten van deze service niet falen indien er geen container genaamd 'redis-container' bestaat (en het proberen verwijderen ervan een fout teruggeeft). Het onderdeel 'ExecStart' is het effectieve commando dat de service zal uitvoeren en zich aan zal binden. In dit geval is dat onze redis container. We geven bij het opstarten via unit files geen flag '-d', detached mee omdat we willen dat de container zich bind aan de service waardoor we de container kunnen monitoren via de aangemaakte service. Hiervoor geven we ook de '--rm' tag mee, deze zorgt ervoor dat bij het stoppen van de container hij ook direct verwijdert wordt. Dit gebeurt jammer genoeg enkel bij een correcte terminatie van de container. Dus indien de docker service crasht zal deze niet verwijderd worden maar wel gestopt.

\section{Managen met tools}

Een derde optie die bestaat is het managen van de containers met tools. In dit onderdeel wordt er niet uitgebreid beschrijven hoe elke tool werkt en hoe die moet gebruikt worden, dit moet geval per geval bekeken worden als dit een meerwaarde betekend voor het project met docker. In dit hoofdstuk wordt er geprobeert een voorbeeld te geven hoe tools zouden kunnen gebruiken worden in een docker omgeving. De tool waarover wat meer uitleg zal gegeven worden is Vagrant. Deze tool is enkel een voorbeeld en is zeker niet de enige tool. Vagrant is een tool ontwikkeld voor het gebruik in ontwikkelomgevingen en wordt door Hashicorp ontwikkeld. Daarnaast hebben ze ook Nomad, een nieuwe tool die in moet staan voor de deployment van containers. Andere bedrijven hebben andere tools en zoals hierboven vermeld moet voor elke use case gekeken worden hoe tools eventueel kunnen geïntegreerd worden in het ontwikkelproces met docker. Het voorbeeld van Vagrant dient om een idee de vormen van hoe tools gebruikt kunnen worden. Docker bied zelf ook een tool aan voor het managen van docker containers namelijk Docker Compose. Een voorbeeld van een configuratiebestand voor Vagrant en Docker Compose worden hieronder weergegeven. Beide zetten een docker container die redis noemt, op van een image genaamd redis.

\begin{lstlisting}[language=bash, style=configstyle]
Vagrant.configure("2") do |config|
	config.vm.provision "docker" do |d|
		d.image = "redis"
		d.name = "redis"
	end
end
\end{lstlisting}
Dit is een voorbeeld van een Vagrant bestand voor docker. Deze kan opgestart worden met het commando 'vagrant up'


\begin{lstlisting}[language=bash, style=configstyle]
#redis image
redis:
	image: redis
\end{lstlisting}
Dit is een voorbeeld van een Docker-compose bestand voor docker. Deze kan opgestart worden met het commando 'docker-compose up'


\chapter{Logging en monitoring}

We kunnen onze Docker containers zo goed beveiligen en managen als we willen, het is bijna nooit de vraag zal er iets verkeerd gaan? Maar eerder de vraag wanneer iets verkeerd zal gaan. Om te kunnen weten wanneer er iets verkeerd gaat en wat er precies verkeerd gaat is het monitoren van onze containers en loggen van wat er gebeurd zeer belangerijk. Daarom heb ik besloten van kort om kort uit te leggen wat een mogelijke manier is om onze containers te monitoren en hoe we onze containers kunnen loggen. De tools die ik hier zal uitleggen zijn zeker niet de enige die hiervoor gebruikt kunnen worden. De focus ligt hier eerder op wat te monitoren en loggen en hoe we dit persistent kunnen doen dan op hoe we dit doen. We zouden kunnen zelf alles opzetten om te monitoren en loggen maar waarom zoveel tijd steken in hoe alles te loggen terwijl we kunnen bezig zijn met het belangerijker aspect namelijk wat te monitoren en loggen \quote{Patrick Debois}. 

Hierbij hebben we net als bij het managen van onze containers weer een aantal opties. 
\chapter{Conclusie}
\label{ch:conclusie}

% TODO: Trek een duidelijke conclusie, in de vorm van een antwoord op de
% onderzoeksvra(a)g(en). Reflecteer kritisch over het resultaat. Zijn er
% zaken die nog niet duidelijk zijn? Heeft het ondezoek geleid tot nieuwe
% vragen die uitnodigen tot verder onderzoek?

Er zijn sinds Docker uitgekomen is heel wat kritieken geweest over de werking van Docker. Zowel gerechtvaardigde als ongerechtvaardigde. Docker biedt ons een tool die flexibeler en complexer is dan het gebruik van een gewone virtuele machine. Er zijn hier wat foute opvattingen over Docker geweest in verband met vooral security maar ook op vlak van managebility. Er werden verwachtingen gesteld dat er met Docker niet meer extra moest beveiligd worden, dat op welke manier we onze containers willen gebruiken, ze automatisch veilig en managable zijn. Op het moment dat wat we normaal op een virtuele machine zouden draaien, gaan containerizen moeten we de beveiliging gewoon op een andere plaats voorzien. Als we gebruik maken van virtuele machines moet er nog altijd gebruik gemaakt worden van beveiliging van het netwerk van de virtuele machines. Daarnaast werd ook elke virtuele machine op een andere manier beveiligd. Deze best practices die we vroeger toepasten op de virtuele machines, mogen we nu niet vergeten bij het gebruik van Docker. Elke Docker container moet nog altijd beveiligd worden naargelang wat we willen uitvoeren op de container. Als we Docker gebruiken maken we gebruik van een flexibele, complexe en krachtige tool. zoals Voltaire zei, "With great power comes great responibility". Docker is een relatief nieuwe tool waarvoor een actieve community is die problemen die opkomen zo snel mogelijk oplossen. Als we onze containers op een virtuele machine uitvoeren en zowel de containers als virtuele machine correct beveiligen en up to date houden, is er zeker een veilige plaats voor Docker in een serveromgeving.

Door de extra flexibiliteit die Docker aanbied, wordt het moelijker om onze containers te managen. Er moet kunnen afgestapt worden van het minutieus controleren van containers. De containers moeten zo gemaakt worden dat er zonder veel menselijke tussenkomst langdurig gebruik van gemaakt kan worden. Er kunnen tools voor het managen van Docker containers zowel voor productie als development gebruikt worden. Bij het gebruik van Docker containers wordt logging en monitoring heel wat belangrijker. Wanneer we onze containers meer willen "loslaten" moet er een systeem bestaan die ervoor instaat ons te waarschuwen als er iets misloopt met onze containers.

Docker bied ons een extra laag van virtualisatie aan. Hierdoor mogen er niet dezelfde eisen stellen aan Docker als dat we doen bij virtuele machines. Beide hebben hun plaats op de markt, en hebben de bedoeling om met elkaar samen te werken en niet elkaar vervangen. Hierbij moeten we rekening houden als we wat vroeger op verschillende virtuele machines zou uitgevoerd worden nu laten uitvoeren op één virtuele machine met Docker containers. We een deel van de controle over deze machines verliezen, maar veel meer flexibiliteit verkrijgen. Daarom mag er niet vergeten worden om de veiligheids elementen die vroeger bestonden om de aparte virtuele machines te beveiligen ook worden meegebracht naar het nieuwe systeem met Docker containers.

\nocite{*}

\bibliographystyle{apa}
\bibliography{bachproef}

%%---------- Back matter -------------------------------------------------

\listoffigures
\listoftables

\end{document}

%moeilijke woorden: linux kernel capabilities
